\chapter{Frontend Auth State (AuthContext)}

\section{Why a Context Instead of Prop Drilling}

Login state is needed across scattered components --- navbar, protected
routes, profile pages. Prop drilling would thread it through components
that don't care about it. React's Context API lets any component subscribe
directly instead.

\begin{diagrambox}[AuthProvider Tree]
\begin{verbatim}
              <AuthProvider>
      -----------------------------------
  <Navbar>   <Dashboard>   <Profile>  <ProtectedRoute>
\end{verbatim}
\end{diagrambox}

\section{AuthContext Implementation}

\begin{lstlisting}[language=JavaScript]
// frontend/src/context/AuthContext.jsx
import { createContext, useContext, useState, useEffect } from "react";
import api from "../services/api";

const AuthContext = createContext(null);

export function AuthProvider({ children }) {
  const [currentUser, setCurrentUser] = useState(null);
  const [loading, setLoading] = useState(true);

  // Restore session on app load by asking the backend who the
  // current cookie/token belongs to, if anyone.
  useEffect(() => {
    const checkAuth = async () => {
      try {
        const { data } = await api.get("/auth/me");
        setCurrentUser(data.user);
      } catch {
        setCurrentUser(null);
      } finally {
        setLoading(false);
      }
    };
    checkAuth();
  }, []);

  const login = async (credentials) => {
    const { data } = await api.post("/auth/login", credentials);
    setCurrentUser(data.user);
    return data.user;
  };

  const logout = async () => {
    await api.post("/auth/logout");
    setCurrentUser(null);
  };

  const value = { currentUser, loading, login, logout };

  return <AuthContext.Provider value={value}>{children}</AuthContext.Provider>;
}

export function useAuth() {
  const context = useContext(AuthContext);
  if (!context) {
    throw new Error("useAuth must be used within an AuthProvider");
  }
  return context;
}
\end{lstlisting}

\begin{notebox}[NOTE]
\texttt{loading} lets \texttt{ProtectedRoute} distinguish ``not known yet''
from ``checked, not logged in,'' avoiding a flash-redirect to
\texttt{/login} on every refresh.
\end{notebox}

\section{Wiring the Provider in main.jsx}

\begin{lstlisting}[language=JavaScript]
// frontend/src/main.jsx
import React from "react";
import ReactDOM from "react-dom/client";
import { BrowserRouter } from "react-router-dom";
import App from "./App";
import { AuthProvider } from "./context/AuthContext";

ReactDOM.createRoot(document.getElementById("root")).render(
  <React.StrictMode>
    <BrowserRouter>
      <AuthProvider>
        <App />
      </AuthProvider>
    </BrowserRouter>
  </React.StrictMode>
);
\end{lstlisting}

\section{Consuming the Context: Navbar}

\begin{lstlisting}[language=JavaScript]
// frontend/src/components/Navbar.jsx
import { useNavigate } from "react-router-dom";
import { useAuth } from "../context/AuthContext";

function Navbar() {
  const { currentUser, logout } = useAuth();
  const navigate = useNavigate();

  const handleLogout = async () => {
    await logout();
    navigate("/login");
  };

  return (
    <nav>
      <span className="brand">Task Manager</span>
      {currentUser ? (
        <div className="navbar-user">
          <span>Hi, {currentUser.name}</span>
          <button onClick={handleLogout}>Logout</button>
        </div>
      ) : (
        <a href="/login">Login</a>
      )}
    </nav>
  );
}

export default Navbar;
\end{lstlisting}

\begin{tipbox}[TIP]
Every component that needs auth state calls \texttt{useAuth()} rather than
importing \texttt{AuthContext} directly, keeping the ``used outside
provider'' guard in one place.
\end{tipbox}
